conserve de joli nom tout cousu d'or de « Fontaourio ».
Car il signifie tout bonnement Fontaine et Source d'Or
(Fons-aurens : Font-aurie) ou source d'or. Cela s'explique
probablement parce que l'eau y coule en abondance. C'est
le cas d'ailleurs d'un hameau qui porte le nom de
« Fontaurie » à quelques kilomètres de Damiatte : et là les
eaux abondent et c'est la fortune du pays...

\par

Ce début de XVII\ieme{} siècle se ressent des temps troublés
qui l'ont précédé. Ce n'est guère favorable à la moralité.
Le brigandage sévit encore. Déjà au XVI\ieme{} siècle les
seigneurs de Paulin et les « Cadets de Padiès » s'étaient
acquis une fâcheuse réputation. Et dès le début de ce
XVII\ieme{} siècle, des actes de banditisme défrayent la
chronique. Dans la nuit du 13 au 14 juillet 1613,
Sébastien de Corbière de Peyrebrune est assailli dans
son château de Saint André. Il ne doit son salut qu'à la
valeur et au courage de domestiques décidés. En 1614,
le même manoir est assiégé et le « Toulouse-Lautrec
Saint Germier » qui s'y est réfugié, n'est délivré que par
les troupes du diocèse d'Albi. En juillet 1616, on jase fort
aux alentours de Paulinet. Le vicomte de Paulin, surpris
en flagrant adultère, a été assassiné par le sieur de Reynes.
Le mari trompé assassine encore le frère de Madame de
Reynès... C'est du beau !

\par

C'est dans cette atmosphère de délinquance sociale que
tout naturellement va éclore une nouvelle épopée
sanglante avec la guerre que patronne le duc de Rohan.